\section{Related Work}
\textbf{Traffic Signal Control.} Early TSC methods such as \textit{FixedTime} rely on predefined cycle lengths and manually designed phase splits for each intersection. Although these approaches are simple and stable in practice, they adapt poorly to time-varying traffic demand. \textit{MaxPressure}\cite{MERCADER2020275} improves adaptability by selecting phases according to the queue imbalance between incoming and outgoing directions, making it more responsive than fixed plans. However, both \textit{FixedTime} and \textit{MaxPressure} remain constrained by handcrafted rules or preset logic, which limits their flexibility under complex and highly dynamic traffic scenarios.

Recent progress in reinforcement learning has led to a series of RL-based TSC methods. \textit{CoLight}\cite{10.1145/3357384.3357902} models network-level cooperation among intersections through graph attention, enabling each signal to coordinate with surrounding intersections instead of acting independently. \textit{CoSLight}\cite{ruan2024coslightcooptimizingcollaboratorselection} further extends this idea by jointly optimizing signal control and collaborator selection, so that each intersection can dynamically decide which neighboring intersections are most useful for coordination under current traffic conditions. \textit{MPLight}\cite{Chen_Wei_Xu_Zheng_Yang_Xiong_Xu_Li_2020} is a decentralized and pressure-aware RL approach designed for large-scale road networks, using parameter sharing to improve scalability and training efficiency. Although these methods achieve strong performance in per city simulation, RL-based controllers face challenges in generalization and interpretability, hence deployment robustness in real-world environments.

To address these limitations, recent studies have begun exploring the integration of large language models (LLMs) into TSC. \textit{LLM-Assisted Light}\cite{wang2024llmassistedlightleveraging} places an LLM at the center of control and augments it with perception and decision-support tools so that it can combine live traffic observations with established TSC methods and remain robust under emergency, roadblock, and sensor-outage scenarios. \textit{iLLM-TSC}\cite{pang2024illmtscintegrationreinforcement} instead adopts a hybrid design in which an RL controller first proposes an action and an LLM then verifies or adjusts that action when observations are degraded or rare events such as packet loss and emergency vehicles fall outside the original RL formulation. \textit{LATS}\cite{zhang2026latslargelanguagemodelassisted} further explores a teacher-student setting in which an LLM teacher provides rich semantic features during multi-agent RL training, while a smaller student controller learns to emulate those features and is deployed without the LLM at inference time. \textit{LLMLight}\cite{lai2024llmlightlargelanguagemodels} moves closer to end-to-end LLM control by adapting its LightGPT backbone through imitation fine-tuning and critic-guided policy refinement so that the model can make direct signal decisions with improved interpretability and cross-scenario performance. \textit{Traffic-R1}\cite{zou2025trafficr1reinforcedllmsbring} extends this trend with a 3B RL-finetuned foundation model that combines expert-aligned offline training, online self-exploration, and asynchronous multi-agent communication to achieve strong zero-shot and out-of-distribution TSC performance. Compared with end-to-end LLM finetuning, LLM+RL hybrids may remain inferior because they still inherit the ceiling and failure modes of the underlying RL controller. Existing end-to-end finetuning methods also have notable limitations: In \textit{LLMLight}, the critic-guided policy refinement objective is weaker than modern policy-gradient methods because the reward/advantage magnitude does not directly participate in loss computation. \textit{Traffic-R1}, meanwhile, relies on substantial human effort to label the offline finetuning data, and its fixed-horizon cumulative queue-based reward assigns the same coarse mean reward to all steps within the horizon.

\textbf{RL finetuning of LLMs.} Large language models are commonly adapted to downstream tasks through instruction finetuning or preference alignment, and instruction finetuning has become a popular baseline because it is simple, stable, and effective at teaching a model to follow task-specific formats and instructions \cite{chung2022scalinginstructionfinetuned,ouyang2022traininglanguagemodels}. By contrast, RL-based finetuning has traditionally been less common in downstream applications because RLHF-style pipelines usually require supervised warm-starting, reward or preference modeling, online rollout collection, and carefully stabilized policy optimization, all of which increase engineering complexity and computational cost \cite{ouyang2022traininglanguagemodels,Sheng_2025}. Recent progress nevertheless shows that reinforcement learning can elicit reasoning capabilities that are difficult to obtain by supervised imitation alone: \textit{DeepSeekMath} introduced Group Relative Policy Optimization (GRPO), which removes the separate value network and replaces it with group-relative advantages to reduce memory overhead while retaining PPO-style policy updates \cite{shao2024deepseekmath}, and \textit{DeepSeek-R1} further demonstrated that large-scale self-exploration can substantially improve reasoning behavior in LLMs \cite{guo2025deepseekr1}. These developments make RL finetuning increasingly attractive for domain-specific agents, since it can reduce dependence on large manually curated datasets, optimize directly for long-horizon behavioral objectives instead of next-token imitation, and preserve useful pretrained abilities through KL-constrained updates around a reference policy \cite{ouyang2022traininglanguagemodels,shao2024deepseekmath,guo2025deepseekr1}.
